% ===================================================================
%  TAFEL Nr. 7 — Das Dorf im Winter. --- Der Bauernhof im Frühling.
%  Traduction 2026 d'apres texto/io, controlee sur texto/fr et
%  texto/en. Consigne : voir texto/de/10-tafel-01.tex.
%
%  ICI LA CONSIGNE DE 2026 TROUVE SA LIMITE, ET IL FAUT L'ECRIRE.
%
%  La regle dit : on modernise la LANGUE, jamais les CHOSES. Cette
%  page nomme deux personnes qui n'existent plus — le Knecht
%  de l'alinea 5 et la Magd de l'alinea 7, le valet et la
%  servante de ferme. Aucun Allemand vivant n'appelle plus personne
%  ainsi ; mais aucun Allemand vivant n'emploie non plus quelqu'un a
%  ces conditions-la. LE MOT N'A PAS VIEILLI PLUS VITE QUE LA CHOSE,
%  il est mort avec elle, et le remplacer par « Angestellter »
%  decrirait un autre monde. On garde donc Knecht et Magd, et l'on
%  ecrit ici pourquoi, pour que ce ne soit pas pris pour un
%  archaisme.
%
%  ET LA PLAISANTERIE DU CORDONNIER, ALINEA 6, CHANGE DE NATURE EN
%  PASSANT EN ALLEMAND. Le vieux Jacob dit : « En ville j'etais
%  cordonnier et je n'avais pas de clients. Ici je suis savetier et
%  j'ai de l'ouvrage. » En francais ce sont DEUX METIERS : le
%  cordonnier fait, le savetier repare. En allemand ce sont DEUX
%  FORMES DU MEME MOT — « Schuhmacher », le faiseur-de-souliers, et
%  « Schuster », qui n'est que sa contraction, du moyen haut-allemand
%  « schuochsūtære », le coud-souliers. La plaisanterie tient encore,
%  mais elle a change de charniere : elle ne porte plus sur le metier,
%  elle porte sur le REGISTRE. Le meme homme, nomme long et court.
%  C'est l'inverse du tableau 5, ou l'allemand avait cinq mots quand
%  l'ido n'en voulait qu'un : ici il n'en a qu'un quand le francais en
%  a deux.
%
%  LA SECONDE SCENE MONTRE LE DERNIER ENDROIT OU UNE LANGUE GARDE
%  TROIS MOTS POUR UNE BETE, et la planche les demande presque tous :
%     Hengst / Stute / Fohlen        l'etalon, la jument, le poulain
%     Stier  / Kuh   / Kalb          le taureau, la vache, le veau
%     Widder / Schaf / Lamm          le belier, la brebis, l'agneau
%     Bock   / Ziege / Zicklein      le bouc, la chevre, le chevreau
%     Eber   / Sau   / Ferkel        le verrat, la truie, le porcelet
%     Hahn   / Henne / Küken         le coq, la poule, le poussin
%     Erpel  / Ente  / Entlein       le canard, la cane, le caneton
%     Ganter / Gans  / Gössel        le jars, l'oie, l'oison
%  Huit series completes, et aucune n'est batie sur la meme racine :
%  ce ne sont pas des derives, ce sont vingt-quatre mots differents.
%  Aucun n'est emprunte. La ferme est la ou une langue depense le plus
%  de mots pour le moins de choses — parce que c'est la qu'on a eu
%  besoin de distinguer, et pendant le plus longtemps.
%
%  Le tableau 2 avait trouve le corps intact et le tableau 5 les
%  metiers intacts. Le betail est le troisieme de ces ilots, et c'est
%  le plus dense des trois.
% ===================================================================

%%K t07-noto-1 noto
\VUnotes{143.2mm}{%
(*) Zu den Einzelheiten einer Mahlzeit siehe die Tafeln 6 und 13.%
}

%%K t07-apar-1 apar
\VUsaut{4.0mm}
\VUcentre{13.2pt}{90}{ZWEITE REIHE}
\VUsaut{3.0mm}
\VUfilet{16mm}
\VUsaut{5.6mm}
\VUtitre{15.0pt}{60}{TAFEL Nr. 7}
\VUsaut{3.0mm}

%%K t07-tit-1 sub
\VUcentre{12.0pt}{0}{\textit{Erste Szene.}}
\VUsaut{3.0mm}

%%K t07-tit-2 sub
\VUpk{10.2pt}{40}{Das Dorf im Winter.}
\VUsaut{4.0mm}

%%K t07-c1-01-1 p
1. --- In den Neujahrsferien bin ich mit meinem Vater nach Neustadt gefahren, in
ein kleines Dorf, wo er etwas zu erledigen hatte.

%%K t07-c1-02-1 p
2. --- Wir nahmen die \VUgras{Postkutsche} \textsuperscript{(11)}, die zwischen der
Stadt und dem Dorf verkehrt; es war aber sehr kalt, und die Fahrt in einem so
unbequemen Wagen war alles andere als angenehm.

%%K t07-c1-03-1 p
3. --- Wir waren also wirklich froh, als der \VUgras{Kutscher} seinen Wagen auf dem
Platz vor dem \VUgras{Gasthaus} \textsuperscript{(2)} \textit{Zum Weißen Bären}
anhielt. Der \VUgras{Wirt} \textsuperscript{(4)} wartete auf der Türschwelle auf uns
und rauchte in aller Ruhe seine \VUgras{Pfeife}.

%%K t07-c1-03-2 p
Er bat uns herein, und ich ließ mich nicht zweimal bitten und setzte mich vor das
Rebholzfeuer, das im Kamin knisterte. Danach war mir der beißende
\VUgras{Nordwind} einerlei, der das alte \VUgras{Schild} \textsuperscript{(3)} knarren
ließ und den \VUgras{Rauch} \textsuperscript{(1)} aus dem Schornstein in schwarzen
Schwaden davontrug.

%%K t07-c1-04-1 p
4. --- Es war Essenszeit. Ohne länger zu warten, sprachen wir dem einfachen, aber
schmackhaften Essen kräftig zu, das man uns auftrug (*).

%%K t07-tit-3 sub
\VUpk{10.2pt}{40}{Ein paar Besuche.}
\VUsaut{3.0mm}

%%K t07-c1-05-1 p
5. --- Nach dem Essen ging ich mit meinem Vater, der Versicherungsvertreter ist, zu
seinen Kunden. Zuerst besuchten wir den \VUgras{Schmied} \textsuperscript{(5)}, der
neben dem Gasthaus wohnt. Er beschlug gerade ein Pferd. Ich bewunderte die Kraft
seines Gehilfen, der den \VUgras{Huf} des Pferdes fest gegen seine Lederschürze
hielt, während der Schmied das \VUgras{Eisen} \textsuperscript{(6)} mit schweren
Hammerschlägen befestigte.

%%K t07-c1-06-1 p
6. --- Dann gingen wir zum \VUgras{Schuhmacher} \textsuperscript{(7)} des Dorfes, dem
alten Jakob. Obwohl er einen prächtigen \VUgras{Stiefel} als Schild hat, begnügte er
sich damit, ein altes Paar abgelaufener Schuhe neu zu besohlen.

%%K t07-c1-06-2 p
Der Alte sagte uns lachend: „In der Stadt war ich ,Schuhmacher‘ und hatte keine
Kunden. Hier bin ich ,Schuster‘ und habe Arbeit genug.“

%%K t07-c1-07-1 p
7. --- Er erzählte uns von seinem Nachbarn, dem \VUgras{Barbier}
\textsuperscript{(8)}, mit dem alle Kunden unzufrieden sind. Seine \VUgras{Rasiermesser}
sind immer schartig und seine \VUgras{Schere} schneidet nie richtig. Aber weil er
der einzige Barbier im Dorf ist, bleibt den Leuten natürlich nichts übrig, als sich
von ihm rasieren und die Haare schneiden zu lassen. Obendrein ist er geizig und
unfreundlich.

%%K t07-c1-08-1 p
8. --- Zum Glück sind die anderen \VUgras{Nachbarn} nicht wie er. Der
\VUgras{Wagner} \textsuperscript{(9)} zum Beispiel ist ein freundlicher Mann,
fleißig und gefällig. Man gibt ihm gern einen Wagen zu reparieren. Seine Arbeit ist
immer sauber abgeliefert.

%%K t07-c1-09-1 p
9. --- Auch über den \VUgras{Sattler} \textsuperscript{(10)}, dem er manchmal beim
Nähen des \VUgras{Geschirrs} hilft, wenn die Arbeit sich häuft, hatte der alte Jakob
nichts zu klagen.

%%K t07-c1-09-2 p
Mein Vater fragte ihn nach seiner Familie. „Allen geht es gut. Meine Enkelin ist in
der \VUgras{Schule} \textsuperscript{(12)}. Ihre \VUgras{Lehrerin}
\textsuperscript{(13)} ist sehr zufrieden mit ihr; sie ist eine gute
\VUgras{Schülerin} \textsuperscript{(14)}. Sie ist nicht wie dieser Taugenichts von
Sohn; der denkt nur ans Spielen. Der \VUgras{Lehrer} \textsuperscript{(15)} bringt
ihm nichts bei.“

%%K t07-tit-4 sub
\VUsaut{3.0mm}
\VUpk{10.2pt}{40}{Dorfgeschichten.}
\VUsaut{3.0mm}

%%K t07-c1-10-1 p
10. --- Der Alte erzählte uns dann die Neuigkeiten aus dem Dorf. Eine neue Schule
sollte gebaut werden, weil die im \VUgras{Gemeindehaus} viel zu klein geworden war.
Das würde leicht gehen: man brauchte nur einen Teil des Gartens des
\VUgras{Müllers} zu kaufen. Dem armen Mann käme das sehr gelegen. Seit dem
Kälteeinbruch konnten die \VUgras{Mühlsteine} der \VUgras{Mühle}
\textsuperscript{(16)} den Weizen nicht mehr mahlen, weil das \VUgras{Rad}
\textsuperscript{(17)} im \VUgras{Eis} \textsuperscript{(25)} des \VUgras{Bachs}
\textsuperscript{(23)} festsaß.

%%K t07-c1-11-1 p
11. --- „Wo wir gerade vom Bauen reden: haben Sie unsere neue \VUgras{Kirche}
\textsuperscript{(18)} gesehen? Unter ihrer Schneehaube sieht sie sehr malerisch aus.
Die \VUgras{Krähen} \textsuperscript{(19)}, die ihren alten Glockenturm nicht mehr
wiedererkennen, sitzen traurig auf dem großen Baum des \VUgras{Friedhofs}
\textsuperscript{(20)}.“

%%K t07-c1-12-1 p
12. --- Der Friedhof brachte den Schuhmacher auf die Leute, die mein Vater gekannt
hatte und die im vergangenen Jahr gestorben waren. „Wie viele \VUgras{Gräber}
\textsuperscript{(21)}, lieber Herr, sind unter der \VUgras{Zypresse}
\textsuperscript{(22)} zu den anderen dazugekommen! Unseren alten Notar hat der Tod
geholt. Das ganze Dorf war bei seiner \VUgras{Beerdigung}.“

%%K t07-c1-13-1 p
13. --- „Da ist sein Nachfolger, der auf dem Bach Schlittschuh läuft. Er ist vor
einer Weile hereingekommen, damit ich den Riemen seines \VUgras{Schlittschuhs}
\textsuperscript{(26)} flicke, der gerissen war.“

%%K t07-c1-14-1 p
14. --- „Und da drüben, am Ufer des Bachs, ist der neue \VUgras{Feldhüter}
\textsuperscript{(27)}. Er ist ein ehemaliger Gendarm und der Schrecken der
Dorflausbuben. Sie laufen davon, sobald sie seine \VUgras{Gamaschen}
\textsuperscript{(29)} und seine \VUgras{Dienstmütze} \textsuperscript{(28)}
erblicken.“

%%K t07-c1-15-1 p
15. --- „Ah, da bringt der alte Josef eine \VUgras{Fuhre Reisig}
\textsuperscript{(30)} für den Bäcker. --- Ganz recht, sagte mein Vater, ich erkenne
sein Gespann \VUgras{Maultiere} \textsuperscript{(31)}. Ich muss mit ihm reden, das
trifft sich gut. Auf Wiedersehen, alter Jakob.“

%%K t07-c1-16-1 p
16. --- Gerade als wir hinausgingen, kamen die Dorfkinder aus der Schule und
stürzten sich auf die \VUgras{Eisbahn} \textsuperscript{(33)}, auf die Gefahr hin,
beim \VUgras{Sturz} \textsuperscript{(34)} hinzufallen und sich etwas zu brechen.
Einer von ihnen holte seinen \VUgras{Schlitten} \textsuperscript{(32)} beim Wagner,
setzte einen Freund darauf und lief los.

%%K t07-c1-17-1 p
17. --- Andere stürzten zu einem \VUgras{Schneehaufen} \textsuperscript{(37)} neben
der \VUgras{Brücke} \textsuperscript{(40)} und bewarfen den \VUgras{Schneemann}
\textsuperscript{(39)}, den sie am Tag zuvor gebaut hatten und dem sie einen Hut,
eine Pfeife und eine Schultasche über der Schulter gegeben hatten.

%%K t07-c1-18-1 p
18. --- Der eisige Wind und die Kälte machten ihnen wenig aus. Die meisten hatten
nicht einmal einen \VUgras{Schal} \textsuperscript{(35)}, und um sich nach der
\VUgras{Schneeball} \textsuperscript{(38)}-Schlacht wieder aufzuwärmen, genügte ihnen
eine Handvoll heißer \VUgras{Kastanien} \textsuperscript{(42)}, die sie bei der alten
Jakobine kauften, der \VUgras{Verkäuferin}, die unter ihrem zu dünnen \VUgras{Tuch}
\textsuperscript{(43)} vor Kälte zittert.

%%K t07-tit-5 sub
\VUsaut{3.0mm}
\VUcentre{12.0pt}{0}{\textit{Zweite Szene.}}
\VUsaut{3.0mm}

%%K t07-tit-6 sub
\VUpk{10.2pt}{40}{Der Bauernhof im Frühling.}
\VUsaut{4.0mm}

%%K t07-c2-01-1 p
1. --- Der Winter hat seine Freuden, aber die Rückkehr des \VUgras{Frühlings}
begrüßen wir doch mit großer Freude.

%%K t07-c2-01-2 p
Was für ein fröhliches Bild ein Bauernhof am Abend im Monat Mai abgibt!

%%K t07-c2-02-1 p
2. --- Am Horizont geht die \VUgras{Sonne} \textsuperscript{(1)} langsam unter und
überflutet das \VUgras{Land} \textsuperscript{(3)} mit ihren purpurnen
\VUgras{Strahlen} \textsuperscript{(2)}. Der fleißige \VUgras{Pflüger}
\textsuperscript{(6)} beeilt sich, sein \VUgras{Feld} \textsuperscript{(4)} zu
pflügen.

%%K t07-c2-02-2 p
Mit seinem \VUgras{Pflug} \textsuperscript{(5)}, den ein kräftiges \VUgras{Pferd}
\textsuperscript{(7)} zieht, zieht er die \VUgras{Furche} \textsuperscript{(8)}, in
die der \VUgras{Sämann} \textsuperscript{(9)} mit vollen Händen das \VUgras{Saatgut}
streuen wird, aus dem die goldenen Ähren werden.

%%K t07-c2-03-1 p
3. --- Die \VUgras{Mäher} \textsuperscript{(11)} haben mit ihren Sensen das Gras der
Wiese geschnitten, die Heuer haben das \VUgras{Heu} \textsuperscript{(10)} mit
\VUgras{Heugabeln} \textsuperscript{(12)} und Rechen gewendet und getrocknet, und
jetzt kommen sie heim.

%%K t07-c2-03-2 p
Die einen gehen vor dem schweren Ochsenkarren her und tragen ihr Werkzeug auf der
Schulter. Die anderen, die fauler sind, haben es sich auf dem duftenden Heu
bequem gemacht.

%%K t07-c2-04-1 p
4. --- Im Vorbeigehen grüßen sie den \VUgras{Gärtner} \textsuperscript{(16)}, der im
\VUgras{Obstgarten} \textsuperscript{(13)} damit beschäftigt ist, die
\VUgras{Obstbäume} zu beschneiden und die \VUgras{Birnbäume}
\textsuperscript{(14)} zu veredeln. Die \VUgras{Pflaumenbäume}
\textsuperscript{(15)} stehen in Blüte, und im gut gedüngten \VUgras{Gemüsegarten}
\textsuperscript{(17)} stehen Gemüse, Kohl und Salat schon gut da.

%%K t07-c2-04-2 p
Auf der niedrigen Mauer schlägt ein schöner \VUgras{Pfau} \textsuperscript{(18)} sein
Rad und zeigt seine prächtigen Federn.

%%K t07-tit-7 sub
\VUpk{10.2pt}{40}{Der Hof.}
\VUsaut{3.0mm}

%%K t07-c2-05-1 p
5. --- Die Türen des \VUgras{Stalls} \textsuperscript{(19)} stehen offen, und man
sieht die Raufe und die \VUgras{Streu} \textsuperscript{(23)} der \VUgras{Stute}
\textsuperscript{(21)}, die der \VUgras{Knecht} \textsuperscript{(20)} eben
ausgespannt hat. Das \VUgras{Fohlen} \textsuperscript{(22)}, auf seinen langen Beinen
so komisch, folgt seiner Mutter und macht Bocksprünge.

%%K t07-c2-06-1 p
6. --- Beim \VUgras{Brunnen} \textsuperscript{(29)} gehen die \VUgras{Kühe}
\textsuperscript{(25)} zum hölzernen \VUgras{Trog} \textsuperscript{(26)} trinken, der
ihnen als Tränke dient. Das \VUgras{Kalb} \textsuperscript{(24)} nutzt die Gelegenheit
zum Saugen, und der \VUgras{Ochse} \textsuperscript{(27)}, der den guten Geruch des
Futters wittert, brüllt fröhlich und hebt stolz den Kopf mit den kräftigen
\VUgras{Hörnern} \textsuperscript{(28)}.

%%K t07-c2-07-1 p
7. --- Alle sind bei der Arbeit. Die \VUgras{Magd} \textsuperscript{(32)} hält das
\VUgras{Seil} \textsuperscript{(31)} des Brunnens. Sie zieht Wasser in einem
\VUgras{Eimer} \textsuperscript{(30)}, um den Tieren zu trinken zu geben. Neben der
\VUgras{Scheune} \textsuperscript{(33)} rafft ein Mann das lose Stroh zusammen, und
vor der Tür des \VUgras{Bauernhauses} \textsuperscript{(34)} spinnt eine alte
\VUgras{Spinnerin} \textsuperscript{(35)} mit ihrem Spinnrad und ihrem
\VUgras{Rocken} den \VUgras{Flachs} oder den \VUgras{Hanf} wie in alten Zeiten.

%%K t07-tit-8 sub
\VUsaut{3.0mm}
\VUpk{10.2pt}{40}{Der Bauer.}
\VUsaut{3.0mm}

%%K t07-c2-08-1 p
8. --- Der \VUgras{Bauer} \textsuperscript{(36)} leitet all diese Arbeit und gibt
seinen Leuten die Anweisungen. Er sagt ihnen, sie sollen die \VUgras{Egge}
\textsuperscript{(40)} und die \VUgras{Spitzhacke} \textsuperscript{(39)}
unterstellen, die neben der \VUgras{Hütte} \textsuperscript{(38)} des \VUgras{Hundes}
\textsuperscript{(37)} liegen geblieben sind.

%%K t07-c2-09-1 p
9. --- Sein Rufen scheucht eine \VUgras{Taube} \textsuperscript{(41)} auf, die in
vollem Flug zum \VUgras{Taubenschlag} \textsuperscript{(42)} flieht, und die
\VUgras{Hühner} \textsuperscript{(43)}, die in den \VUgras{Hühnerstall}
\textsuperscript{(44)} zurückeilen.

%%K t07-c2-10-1 p
10. --- Vor dem \VUgras{Schafstall} \textsuperscript{(48)} bringt der alte
\VUgras{Schäfer} \textsuperscript{(45)} seine \VUgras{Herde} \textsuperscript{(49)}
zurück. Mit seinem \VUgras{Hirtenstab} \textsuperscript{(46)} in der Hand, der ihn so
wenig verlässt wie sein verschossener \VUgras{Kittel} \textsuperscript{(47)}, führt
er \VUgras{Widder} \textsuperscript{(50)}, \VUgras{Mutterschafe}
\textsuperscript{(53)}, \VUgras{Lämmer} \textsuperscript{(54)}, \VUgras{Ziegen}
\textsuperscript{(51)} und \VUgras{Zicklein} \textsuperscript{(52)} in den Pferch.
Sein treuer \VUgras{Hund} \textsuperscript{(55)} Argus kommt zuletzt und treibt die
Nachzügler mit seinem Bellen an.

%%K t07-c2-11-1 p
11. --- All dieser Lärm und Trubel stört die Ruhe der guten \VUgras{Milchkuh}
\textsuperscript{(56)} nicht, die ihre \VUgras{Glocke} \textsuperscript{(57)} läuten
lässt, während sie vor der Tür des \VUgras{Kuhstalls} \textsuperscript{(58)} gemolken
wird.

%%K t07-c2-12-1 p
12. --- Sie scheint zu verstehen, dass sie ihre Milch geben muss, damit die
\VUgras{Bäuerin} \textsuperscript{(60)} im \VUgras{Butterfass}
\textsuperscript{(61)} \VUgras{Butter} machen kann oder die ausgezeichneten
\VUgras{Käse} \textsuperscript{(64)}, die auf dem Brett über dem \VUgras{Zuber}
\textsuperscript{(63)} abtropfen.

%%K t07-c2-13-1 p
13. --- Die ganze \VUgras{Milchkammer} \textsuperscript{(59)} wird vom
\VUgras{Hofknecht} \textsuperscript{(65)} sehr sauber gehalten, der eben die
\VUgras{Milchkannen} \textsuperscript{(62)} im großen Eimer gewaschen hat, den er
trägt.

%%K t07-tit-9 sub
\VUsaut{3.0mm}
\VUpk{10.2pt}{40}{Der Geflügelhof.}
\VUsaut{3.0mm}

%%K t07-c2-14-1 p
14. --- In seinen schweren Holzschuhen geht er ziemlich ungeschickt und verscheucht
damit den \VUgras{Truthahn} \textsuperscript{(68)}, die \VUgras{Enten}
\textsuperscript{(66)}, die \VUgras{Erpel} \textsuperscript{(67)} und die
\VUgras{Gänse} \textsuperscript{(69)}, die ihre \VUgras{Schnäbel}
\textsuperscript{(70)} weit aufreißen und ein ängstliches Geschrei anstimmen.

%%K t07-c2-14-2 p
Der \VUgras{Hahn} \textsuperscript{(71)}, kriegerischer, hebt seinen roten
\VUgras{Kamm} \textsuperscript{(72)}, schüttelt die Federn seines
\VUgras{Schwanzes} \textsuperscript{(73)} und lässt ein helles „Kikeriki“ hören.

%%K t07-c2-15-1 p
15. --- Eine \VUgras{Glucke} \textsuperscript{(74)} scharrt auf dem
\VUgras{Misthaufen} \textsuperscript{(81)} mit ihren \VUgras{Küken}
\textsuperscript{(75)}. Das
\VUgras{Ferkel} \textsuperscript{(78)} läuft zu seiner Mutter, der riesigen
\VUgras{Sau} \textsuperscript{(77)}, die wir neben dem Schweinestall sehen.

%%K t07-c2-16-1 p
16. --- In ihren Ställen fressen die \VUgras{Kaninchen} \textsuperscript{(79)} gierig
das zarte \VUgras{Gras} \textsuperscript{(80)}, das die Tochter des Bauern ihnen
bringt. Hannchen liebt ihre Kaninchen sehr, und jeden Tag, nachdem sie die frischen
\VUgras{Eier} \textsuperscript{(76)} aus dem Hühnerstall geholt hat, kommt sie ihre
kleinen Freunde besuchen.
\VUsaut{6.0mm}
\VUfilet{22mm}
